%%
%% Chapter 5: Conclusion and Outlook (结论与展望)
%%

\chapter{结论与展望}

\section{研究结论}

本研究通过理论分析和实证研究，得出以下主要结论：

第一，变量A对变量B具有显著的正向影响，这一结论与已有研究结果一致，进一步验证了相关理论。

第二，变量C在变量A与变量B的关系中起到了中介作用，这揭示了两者关系的内在机制。

第三，变量D的调节效应在特定条件下显著，这为边界条件的探讨提供了新的视角。

\section{理论贡献}

本研究的理论贡献主要体现在：丰富了相关领域的理论研究，为后续研究提供了新的分析框架；揭示了变量间关系的内在机制和边界条件，扩展了现有理论的适用范围。

\section{实践启示}

本研究的实践启示包括：为组织管理者提供了基于实证证据的决策参考；为政策制定者提供了改进相关政策的方向和思路。

\section{研究局限与未来展望}

本研究存在以下局限性：样本来源较为单一，可能影响研究结论的外部效度；采用横截面研究设计，难以严格推断变量间的因果关系。

未来研究可以从以下方向展开：扩大样本来源范围，提高研究结论的普适性；采用纵向研究设计，揭示变量间的动态变化关系；引入更多可能的解释变量，完善研究模型。
